\chapter*{CAPÍTULO 1. BASES CONCEPTUALES}
\addcontentsline{toc}{chapter}{CAPÍTULO 1. BASES CONCEPTUALES}
\setcounter{chapter}{1}
\setcounter{section}{0}
\renewcommand{\thesection}{\thechapter.\arabic{section}}
\renewcommand{\thesubsection}{\thesection.\arabic{subsection}}
\renewcommand{\thesubsubsection}{\thesubsection.\arabic{subsubsection}}

\section{Deuda social en el desarrollo de software}

La deuda social es un fenómeno que ha ganado relevancia en los equipos de desarrollo de software ágil debido a su impacto significativo y las consecuencias negativas que conlleva. Entre estas consecuencias se encuentran: la deficiente comunicación entre los miembros del equipo, el desgaste de los desarrolladores por la presión cons-tante de cumplir plazos estrictos, la falta de reconocimiento del esfuerzo individual y la ausencia de procesos de retroalimentación que dificultan la mejora continua. Según Espinosa et al. \cite{caballero2022communitysmells}, la deuda social aborda cómo las decisiones técnicas y sociales impactan los entornos de trabajo y a las personas, generándose a partir de decisiones sociotécnicas deficientes o subóptimas. En otras palabras, a medida que se toman decisiones erróneas de manera constante, los efectos negativos de la deuda social en los equipos de desarrollo de software se vuelven más pronunciados, y mientras estas malas decisiones persistan, sus efectos se intensificarán con el tiempo \cite{caballero2021thesis}.
En relación con la deuda social en el desarrollo de software ágil, existen términos asociados, como los olores comunitarios. Según Tamburri et al. \cite{tamburri2015insights}, los olores comunitarios son “conjuntos de circunstancias organizativas y sociales que tienen relaciones causales implícitas” \cite{tamburri2015insights}. Estos olores comunitarios son señales de problemas en los equipos de desarrollo de software, como conflictos entre miembros, falta de normas o procesos, o una cultura de trabajo tóxica \cite{caballero2022communitysmells}, \cite{caballero2021thesis}. Si estos problemas no se abordan, pueden conducir a la acumulación de deuda social \cite{martini2019agiledebt}.

\section{Comunicación, colaboración y coordinación (3C)}

La comunicación, colaboración y coordinación (3C) son aspectos clave al hablar de deuda social. La ausencia de estos elementos puede generar malestar en los equipos de desarrollo \cite{caballero2022communitysmells}. En cuanto a la comunicación, Saeeda et al. \cite{saeeda2024nontechnicaldebt} afirman que una comunicación efectiva es esencial para el éxito en el desarrollo de software. La adopción de atajos en la comunicación puede llevar a errores y retrasos, como omitir pasos en la recopilación de requisitos del cliente, lo que resulta en un análisis incompleto o incorrecto \cite{saeeda2024nontechnicaldebt}. De manera similar, la falta de colaboración puede tener consecuencias negativas tanto para la empresa como para los empleados. Sin una cultura de colaboración, los equipos pueden trabajar de manera aislada y descoordinada, generando errores, retrasos y menor productividad \cite{saeeda2024nontechnicaldebt}. Por último, la falta de coordinación entre los miembros del equipo puede ralentizar el progreso, aumentar los errores y reducir la productividad \cite{saeeda2024nontechnicaldebt}.

\section{Compromiso organizacional}

El compromiso organizacional es una variable clave en la gestión de la deuda social. Dignan \cite{dignan2020organizationaldebt} sostiene que la falta de compromiso por parte de la organización puede tener un impacto negativo en el desarrollo de software. Una organización ineficaz puede dificultar la comunicación y la colaboración, lo que puede llevar a errores y retrasos. Además, la falta de recursos para la implementación y mantenimiento del software puede afectar la calidad del producto final \cite{dignan2020organizationaldebt}. En situaciones como esta, es necesario un cambio. Sin embargo, como alerta Saeeda et al. \cite{saeeda2024nontechnicaldebt}, es importante prepararse para la resistencia al cambio. Para superarla, es fundamental comunicar los motivos del cambio de manera clara y convincente. Aunque los cambios importantes suelen justificarse por beneficios financieros, es crucial que las personas comprendan cómo estos beneficios se traducen en ventajas tangibles para ellos \cite{saeeda2024nontechnicaldebt}.

\section{Síndrome de Burnout}

El síndrome de Burnout es un problema frecuente en el desarrollo de software ágil y afecta gravemente a las personas. Según Tulili et al. \cite{tulili2022burnout}, el síndrome de Burnout es el agotamiento mental, físico y emocional provocado por factores relacionados con el trabajo. Los factores más comunes incluyen la tensión y el control en el trabajo, la sobrecarga laboral, el conflicto de roles, la especialización del trabajo y las exigencias laborales \cite{tulili2022burnout}. Otros estudios han identificado factores adicionales que contribuyen al Burnout, como la sobrecarga de información, la cascada de datos, largas rachas de contribuciones ininterrumpidas y las tensiones laborales prolongadas derivadas de la pandemia de COVID-19 \cite{tulili2022burnout}. Al ser un problema que afecta directamente a las personas, el síndrome de Burnout puede generar deuda social y, por lo tanto, debe considerarse al gestionar esta deuda.

\section{Goal Question Metric (GQM)}

Es una metodología estructurada para definir y evaluar métricas de software alineadas con los objetivos organizacionales. Desarrollado por Victor R. Basili, Gianluigi Caldiera y H. Dieter Rombach en 1994 \cite{caldiera1994gqm}, este enfoque se basa en tres niveles: definir metas claras y específicas (Goal), formular preguntas que ayuden a evaluar si se están alcanzando las metas (Question) e identificar y recolectar datos que respondan a estas preguntas (Metric). El GQM permite a las organizaciones medir de manera efectiva e interpretar los datos en función de sus objetivos, mejorando así la calidad del software a través de prácticas de medición sistemáticas.

\section{Scrum}

Scrum es un marco de trabajo que ayuda a las personas, equipos y organizaciones a generar valor a través de soluciones adaptativas para problemas complejos; donde su arquitectura en equipos reducidos y autogestionados operan mediante ciclos iterativos —denominados Sprints—, los cuales, tiene una duración máxima de un mes, priorizando la entrega iterativa de valor. En Scrum, tres roles vertebran su estructura: el Scrum Master, que facilita los procesos; el Product Owner, el responsable del valor; y el Equipo de Desarrollo, ejecutores de tareas. Su núcleo operativo reside en la tríada de transparencia-inspección-adaptación. En eventos estructurados tenemos el Sprint Planning, Daily Scrum, Sprint Review y Sprint Retrospective \cite{schwaber2020scrum}. Cabe subrayar que, a diferencia de otras metodologías, Scrum no prescribe herramientas técnicas específicas; su esencia radica en proveer un esqueleto mínimo que catalice la colaboración y la mejora progresiva \cite{takeuchi1986new}.